\subsubsection{Збіжність чебишовської інтерполяції для диференційовних функцій}

Центральне положення теорії наближень полягає в тому, що швидкість збіжності поліноміальних інтерполянтів за вузлами Чебишева при $n \to \infty$ безпосередньо визначається властивостями гладкості аналізованої функції $f(x)$. Якщо функція володіє обмеженою кількістю похідних, швидкість спадання похибки має чітко виражений алгебраїчний характер.

Зокрема, якщо функція $f$ є диференційовною $\nu$ разів на інтервалі $[-1, 1]$, то швидкість збіжності її чебишовського інтерполянта у рівномірній нормі залежить від поведінки вищих похідних. Класичним прикладом функції обмеженої гладкості є $f(x) = |x|$, яка має розрив першої похідної в точці $x=0$ і демонструє лінійну швидкість спадання похибки $\mathcal{O}(n^{-1})$.

Для строгого аналізу швидкості збіжності у випадках, коли вищі похідні можуть мати точки розриву або скачки, використовується математичний апарат функцій з обмеженою повною варіацією (bounded variation). Якщо $\nu$-та похідна функції має скінченну повну варіацію $V$ на відрізку $[-1, 1]$, це дозволяє отримати точну аналітичну оцінку як для коефіцієнтів розкладу Чебишева, так і для похибки самого інтерполяційного полінома.

\begin{statement}[Збіжність для функцій з обмеженою варіацією похідної]
	\label{st:convergence_differentiable}
	Нехай функція $f$ та її похідні до порядку $\nu-1$ включно є абсолютно неперервними на відрізку $[-1, 1]$, а похідна порядку $\nu$ ($f^{(\nu)}$) має обмежену повну варіацію $V$. 
	
	Тоді коефіцієнти $a_k$ розкладу цієї функції в ряд за поліномами Чебишева задовольняють нерівність:
	\begin{equation}\label{eq:cheb_coeff_bound}
		|a_k| \le \frac{2V}{\pi k^{\nu+1}}, \quad k \ge 1.
	\end{equation}
	
	Окрім того, для будь-якого $n > \nu$ похибка чебишовського інтерполяційного полінома $p_n$ ступеня $n$ у рівномірній нормі обмежена зверху і задовольняє співвідношення:
	\begin{equation}\label{eq:cheb_interpolant_error}
		\|f - p_n\| \le \frac{4V}{\pi \nu (n-\nu)^\nu}.
	\end{equation}
\end{statement}

Ці оцінки показують, що барицентрична інтерполяція за вузлами Чебишева гарантує стійке та передбачуване спадання похибки зі збільшенням степеня полінома $n$. Алгебраїчна швидкість збіжності повністю визначається глобальними та локальними диференціальними властивостями функції, що дозволяє строго контролювати точність наближення на всьому інтервалі інтерполяції.